\title{
MATEMATIK
}
Chalmers tekniska högskola
Datum: 2020-10-10 kl. 08:30-12.30
Tentamen
Telefonvakt: Hossein Raufi
Telefon: 070 - 4490237
\section*{MVE041 Flervariabelmatematik Z1}
Betygsgränser: 3: 20-29 p, 4: 30-39, 5: 40-50.
Lösningar läggs ut på kursens webbsida första vardagen efter tentamensdagen.
Till samtliga uppgifter skall fullständiga lösningar inlämnas. Endast svar ger inga poäng. Motivera och förklara så väl du kan.
1. (a) Beräkna \(\iint_{D} \sin (x+y) d A\) där \(D\) är området som beskrivs av olikheterna \(y \geq 0\), \(y \leq x \operatorname{och} y \leq 2-x\).
(a) Beräkna \(\iint_{T} x y \sqrt{x^{2}+y^{2}} d A\) där \(T\) är triangeln med hörnen \((0,0),(1,0)\) och \((1,1)\). Använd därefter detta resultat till att beräkna \(\iint_{D} x y \sqrt{x^{2}+y^{2}} d A\) där \(D\) är kvadraten med hörnen \((0,0),(1,0),(0,1)\) och \((1,1)\).
2. (a) Bestäm tangentplanet till ytan \(x^{2} z+z^{2}=2 x y+y^{2}\) i punkten \((0,-1,1)\).
(b) Bestäm \(C \in \mathbb{R}\) så att planet \(2 x+2 y+z=C\) tangerar ytan \(x+y^{2}+z^{4}=1\).
3. Beräkna arbetet som utförs då kraftfältet
\[
\mathbb{F}(x, y)=\left((2 x+x^{2} y) e^{x y}, x^{3} e^{x y}+1\right)
\]
flyttar en partikel längs kurvan \(y=\sqrt{\frac{9}{4}-9 x^{2}}\) från \(x=0\) till \(x=-1 / 2\).
4. Transformera den partiella differentialekvationen
\[
x^{2} \frac{\partial^{2} f}{\partial x^{2}}+2 x y \frac{\partial^{2} f}{\partial x \partial y}+y^{2} \frac{\partial^{2} f}{\partial y^{2}}=x y, \quad x>0
\]
genom att göra variabelbytet
\[
\left\{\begin{array}{l}
u=x \\
v=\frac{x}{y}
\end{array}\right.
\]
Obs! Du behöver inte lösa differentialekvationen, bara transformera den från \((x, y)\) till \((u, v)\)-koordinator.
5. Använd Lagrange-multiplikatorer till att bestämma den största möjliga volymen för ett rätblock vars totala begränsningsyta är \(2 \mathrm{~m}^{2}\).
Var god vänd!

6. Berākna volymen av den kropp som beskrivs av olikheterna \(z \geq 0, z \geq x^{2}+y^{2}-1\), och \(z \leq \sqrt{x^{2}+y^{2}}+1\).
7. Ytan \(z=\sqrt{4-x^{2}-y^{2}}, x \geq 0, y \geq 0\) och \(x^{2}+y^{2} \leq 4\) har densiteten
\[
\rho(x, y, z)=z \sqrt{x^{2}+y^{2}} \mathrm{kg} / \mathrm{m}^{2} .
\]
Berākna ytans massa.
8. Berākna det största värde som linjeintegralen
\[
\oint_{\gamma} y^{3} d x + \left(3 x-x^{3}\right) d y
\]
där \(\gamma\) är en enkel, sluten kurva genomlöpt ett varv moturs, kan anta.
Lycka till!
/Hossein